\chapter{Reference}
\label{ch:reference}

This chapter provides quick-reference material for day-to-day Veilbox
operation, including command cheat sheets, configuration file summaries,
and troubleshooting guides.

\section{Command Cheat Sheet}

\subsection{Build Commands}
\begin{longtable}{|p{5cm}|p{8cm}|}
\hline
\textbf{Command} & \textbf{Description} \\
\hline
./build.sh & Full build (all 14 stages). \\
./build.sh --clean & Clean all build artifacts. \\
./build.sh --stage kernel & Build kernel only (incremental). \\
./build.sh --stage disk & Build disk image only. \\
./build.sh --stage initramfs & Regenerate initramfs file list. \\
MEM=2G ./test.sh & Boot in QEMU (direct kernel). \\
./test.sh --bios & Boot via SeaBIOS + GRUB. \\
./test.sh --check & Auto health check (exit 0 on login). \\
MEM=2G ./test.sh --keep-state & Boot with persistent state disk. \\
\hline
\end{longtable}

\subsection{SSH Commands}
\begin{longtable}{|p{5cm}|p{8cm}|}
\hline
\textbf{Command} & \textbf{Description} \\
\hline
ssh -i output/ssh-test-key root@localhost -p 2222 & Key auth. \\
ssh root@localhost -p 2222 & Password auth (veiladmin). \\
ssh-keygen -R '[localhost]:2222' & Remove old host key. \\
\hline
\end{longtable}

\subsection{Container Commands}
\begin{longtable}{|p{5cm}|p{8cm}|}
\hline
\textbf{Command} & \textbf{Description} \\
\hline
nerdctl images & List pulled images. \\
nerdctl pull <image> & Pull an image. \\
nerdctl run <image> <cmd> & Run a container. \\
nerdctl run -it <image> sh & Interactive shell. \\
nerdctl run -d --name <n> <image> & Detached container. \\
nerdctl ps & List running containers. \\
nerdctl ps -a & List all containers. \\
nerdctl stop <id> & Stop container. \\
nerdctl rm <id> & Remove container. \\
nerdctl exec <id> <cmd> & Exec in running container. \\
nerdctl logs <id> & View container logs. \\
nerdctl stats & Resource usage. \\
nerdctl inspect <id> & Detailed info. \\
\hline
\end{longtable}

\subsection{System Commands}
\begin{longtable}{|p{5cm}|p{8cm}|}
\hline
\textbf{Command} & \textbf{Description} \\
\hline
ps & List processes. \\
ifconfig eth0 & Show IP address. \\
netstat -tlnp & Show listening ports. \\
df -h & Show disk usage. \\
free -m & Show memory usage. \\
mount & Show mounted filesystems. \\
cat /var/log/messages & View system log. \\
poweroff & Shutdown. \\
reboot & Reboot. \\
\hline
\end{longtable}

\section{Configuration File Reference}

\begin{longtable}{|p{4cm}|p{3cm}|p{5cm}|}
\hline
\textbf{File} & \textbf{Format} & \textbf{Purpose} \\
\hline
/etc/inittab & BusyBox init & Spawn getty, run rcS, shutdown. \\
/etc/init.d/rcS & Shell script & System initialization. \\
/etc/passwd & Colon-delimited & User accounts. \\
/etc/shadow & Colon-delimited & Password hashes. \\
/etc/hostname & Plain text & System hostname. \\
/etc/containerd/config.toml & TOML & Container runtime config. \\
/root/.profile & Shell script & User login prompt. \\
/usr/share/udhcpc/default.script & Shell script & DHCP configuration. \\
\hline
\end{longtable}

\section{Kernel Config Key Reference}

\begin{longtable}{|p{5cm}|p{3cm}|p{5cm}|}
\hline
\textbf{Option} & \textbf{Value} & \textbf{Purpose} \\
\hline
CONFIG\_NAMESPACES & y & Container isolation. \\
CONFIG\_CGROUPS & y & Resource limits. \\
CONFIG\_OVERLAY\_FS & y & Image layers. \\
CONFIG\_VETH & y & Container networking. \\
CONFIG\_BRIDGE & y & Network bridge. \\
CONFIG\_VIRTIO\_NET & y & VirtIO NIC. \\
CONFIG\_VIRTIO\_BLK & y & VirtIO disk. \\
CONFIG\_EXT4\_FS & y & Root filesystem. \\
CONFIG\_INITRAMFS\_SOURCE & "initramfs.list" & Embedded rootfs. \\
CONFIG\_EFI\_STUB & y & QEMU -kernel boot. \\
CONFIG\_SERIAL\_8250\_CONSOLE & y & Serial console. \\
CONFIG\_SECCOMP & y & Syscall filtering. \\
CONFIG\_DEVTMPFS & y & Device node creation. \\
\hline
\end{longtable}

\section{BusyBox Applet Reference}

\begin{longtable}{|p{3cm}|p{4cm}|p{5cm}|}
\hline
\textbf{Category} & \textbf{Applets} & \textbf{Notes} \\
\hline
Shell & ash, sh & POSIX shell (no bashisms). \\
File & ls, cp, mv, rm, mkdir, rmdir, cat, chmod, chown, ln, find, grep, sed, awk & Core utils. \\
Process & ps, kill, pgrep, pkill, top & Process management. \\
Network & ifconfig, route, ping, netstat, udhcpc, wget & Networking. \\
Init & init, getty, login, passwd & System init. \\
Editor & vi & Minimal vi. \\
Archive & tar, gzip, gunzip, bzip2, bunzip2, xz, unxz, cpio & Archive tools. \\
Misc & echo, printf, test, true, false, sleep, timeout, env, set, export, read, exit & Shell builtins. \\
\hline
\end{longtable}

\section{Kernel Parameter Reference}

\begin{longtable}{|p{6cm}|p{7cm}|}
\hline
\textbf{Parameter} & \textbf{Format} \\
\hline
console & \texttt{console=device,options} \\
root & \texttt{root=/dev/device} \\
init & \texttt{init=/path/to/init} \\
quiet & (flag, reduces boot messages) \\
debug & (flag, maximum verbosity) \\
single & (flag, single-user mode) \\
panic & \texttt{panic=seconds} \\
veilbox.autologin & (flag, custom Veilbox auto-login) \\
net.ifnames & \texttt{net.ifnames=0} (legacy ethX names) \\
mitigations & \texttt{mitigations=off} (disable CPU vuln fixes) \\
clocksource & \texttt{clocksource=tsc} (use TSC timer) \\
\hline
\end{longtable}

\section{Build Script Environment Variables}

\begin{longtable}{|p{4cm}|p{4cm}|p{5cm}|}
\hline
\textbf{Variable} & \textbf{Default} & \textbf{Purpose} \\
\hline
MEM & 2G & Guest RAM for QEMU \\
OUTPUT & output/ & Build output directory \\
ROOTFS & rootfs/ & Root filesystem source \\
KERNEL\_SRC & linux-7.1.0-rc6/ & Kernel source directory \\
CACHE\_DIR & \$HOME/veilbox-cache & Download cache \\
VERBOSE & 0 & Enable verbose build output \\
SKIP\_DEPS & 0 & Skip dependency check \\
\hline
\end{longtable}

\section{Port Reference}

\begin{longtable}{|p{3cm}|p{3cm}|p{6cm}|}
\hline
\textbf{Port} & \textbf{Protocol} & \textbf{Service} \\
\hline
22 & TCP & SSH (Dropbear) — guest internal \\
2222 & TCP & SSH (QEMU forwarded from host) \\
80 & TCP & HTTP (container web servers) \\
1338 & TCP & Containerd metrics (localhost only) \\
3306 & TCP & MariaDB/MySQL (container) \\
5000 & TCP & Flask/Python apps (container) \\
5201 & TCP & iperf3 (container network test) \\
\hline
\end{longtable}

\section{File System Layout Reference}

\begin{lstlisting}
# Complete filesystem hierarchy
/               # initramfs root (tmpfs)
|------ bin/        # BusyBox applets
|------ sbin/       # System binaries (init, getty)
|------ etc/        # Configuration
|   |------ inittab
|   |------ init.d/rcS
|   |------ passwd
|   |------ shadow
|   |------ hostname
|   |------ resolv.conf
|   `------ dropbear/
|       |------ authorized_keys
|       |------ dropbear_ed25519_host_key
|       `------ dropbear_rsa_host_key
|------ usr/
|   |------ bin/    # containerd, nerdctl (aliased as docker), dropbear, dropbearkey
|   `------ sbin/
|------ lib64/      # Shared libraries (glibc, libseccomp)
|------ proc/       # procfs
|------ sys/        # sysfs
|------ dev/        # devtmpfs
|------ tmp/        # tmpfs
|------ run/        # tmpfs (runtime data)
|------ var/
|   `------ log/
|       `------ messages  # syslog output
|------ mnt/
|   `------ state/  # Persistent state (ext4 or tmpfs fallback)
`------ root/       # Root user home
    `------ .profile
\end{lstlisting}

\section{Quick Commands Reference}

\begin{longtable}{|p{5cm}|p{8cm}|}
\hline
\textbf{Action} & \textbf{Command} \\
\hline
Build everything & \texttt{./build.sh} \\
Clean build & \texttt{./build.sh --clean} \\
Boot in QEMU & \texttt{MEM=2G ./test.sh} \\
Boot with state & \texttt{./test.sh --keep-state} \\
Auto health check & \texttt{./test.sh --check} \\
GRUB BIOS boot & \texttt{./test.sh --bios} \\
SSH connect & \texttt{ssh -i output/ssh-test-key root@localhost -p 2222} \\
Kernel menuconfig & \texttt{cd linux-7.1.0-rc6 \&\& make menuconfig} \\
Check logs & \texttt{cat /var/log/messages} \\
Free memory & \texttt{free -m} \\
Disk usage & \texttt{df -h} \\
List processes & \texttt{ps} \\
Network config & \texttt{ifconfig eth0} \\
Routing table & \texttt{route -n} \\
\hline
\end{longtable}
